\numberlesssection{ВВЕДЕНИЕ}

Современные системы онлайн-доставки еды ежедневно обрабатывают большое количество заказов, координируют работу тысяч ресторанов и курьеров, управляют платежами и отслеживают статус доставок. Основной причиной развития таких систем является повседневное использование мобильных приложений, а также пандемия COVID-19, закрепившая у пользователей привычку к бесконтактным услугам~\cite{ref:covid-19}.

Согласно актуальным статистическим данным, российский рынок готовой еды в 2025 году демонстрирует устойчивую тенденцию к росту. Ключевым фактором развития отрасли является не только рост спроса со стороны потребителей, но и значительное увеличение выручки производителей. При этом рынок сталкивается с рядом вызовов: ростом себестоимости, дефицитом производственных мощностей и необходимостью цифровизации для оптимизации логистических процессов~\cite{ref:growing}.

В основе таких систем лежит база данных, которая должна обеспечивать быстрый доступ к информации, гарантировать целостность финансовых транзакций и хранить историю операций. От качества проектирования базы данных зависит скорость обработки заказов, точность расчётов и возможность масштабирования системы.

Цель работы: разработка базы данных для системы онлайн-заказов еды, а также приложения для доступа к ней.

Для достижения поставленной цели необходимо решить следующие задачи:
\begin{itemize}
    \item провести анализ существующих систем онлайн-заказов еды;
    \item определить требования и ограничения к разрабатываемой базе данных и приложению;
    \item спроектировать базу данных и задать ограничения целостности;
    \item выбрать средства реализации базы данных и приложения;
    \item реализовать спроектированные базу данных и приложение;
    \item провести исследование влияния внешнего кэширования и индексации на время выполнения запросов к базе данных.
\end{itemize}
